% GENERATED FILE. Edit canonical lessons, then run npm run generate:books.
% canonical-source-hash: fnv1a64:8a4551aa76219748
% canonical-lessons: IT-C438-reclamo, IT-C438-reclamare, IT-C438-lamentela, IT-C438-rimediare, IT-C438-risarcire

\chapter{Il reclamo}
\label{ch:it-il-reclamo}
\hlchaptermodality{438}

\begin{chapteropening}
\textbf{By the end of this chapter:} \emph{I can make a complaint and get it put right: the complaint, complaining, a grumble, putting it right and compensation.}

Complete the last lesson of chapter 438: I can make a complaint and get it put right: the complaint, complaining, a grumble, putting it right and compensation.
\end{chapteropening}

\section[il reclamo]{il reclamo — the complaint (formal)}

\label{lesson:IT-C438-reclamo}



\begin{quote}
Before the new one: say the Italian for the mould, then the Italian for to creak.
\end{quote}

\subsection*{You'll want to know: il reclamo}
\textbf{il reclamo} — \textquotedblleft{}the complaint (formal)\textquotedblright{}.

\begin{quote}
Ho scritto un reclamo alla direzione.
\end{quote}

\begin{quote}
\emph{I wrote a complaint to the management.}
\end{quote}

\emph{Fare reclamo} — to lodge a complaint.

\begin{cousinweb}
\emph{reclamo} is built on \emph{reclamare}, a learned borrowing from Latin \emph{reclāmāre} 'to cry out against' (EtymDB), from \emph{clāmāre} 'to shout'.
\end{cousinweb}

\begin{grammarlens}[title={What you've built}]
One more word for telling, explaining and describing at B1.
\end{grammarlens}

\subsection*{Your turn}
\begin{itemize}
\raggedright
  \item \emph{Say it:} \emph{il reclamo}
  \item \emph{Say it:} \emph{il reclamo}, once more
  \item \emph{Say it:} a sentence of your own with \emph{il reclamo}
  \item \emph{Recall:} say the Italian for the mould, then the Italian for to creak, then say \emph{il reclamo} again
\end{itemize}

\begin{tcolorbox}[breakable,colback=teal!4,colframe=teal!35!black,title={Before you move on}]
What does il reclamo mean? (\textquotedblleft{}The complaint (formal)\textquotedblright{}.) Say it once more.
\end{tcolorbox}
\section[reclamare]{reclamare — to complain, to protest; to claim}

\label{lesson:IT-C438-reclamare}



\begin{quote}
Before the new one: say the Italian for to creak, then the Italian for the complaint.
\end{quote}

\subsection*{You'll want to know: reclamare}
\textbf{reclamare} — \textquotedblleft{}to complain, to protest; to claim\textquotedblright{}.

\begin{quote}
I clienti reclamano per il rumore.
\end{quote}

\begin{quote}
\emph{The guests complain about the noise.}
\end{quote}

\emph{Reclamare un diritto} — to claim a right.

\begin{cousinweb}
\emph{reclamare} is a learned borrowing from Latin \emph{reclāmāre} 'to cry out against' (EtymDB), from \emph{clāmāre} 'to shout'.
\end{cousinweb}

\begin{grammarlens}[title={What you've built}]
One more word for telling, explaining and describing at B1.
\end{grammarlens}

\subsection*{Your turn}
\begin{itemize}
\raggedright
  \item \emph{Say it:} \emph{reclamare}
  \item \emph{Say it:} \emph{reclamare}, once more
  \item \emph{Say it:} a sentence of your own with \emph{reclamare}
  \item \emph{Recall:} say the Italian for to creak, then the Italian for the complaint, then say \emph{reclamare} again
\end{itemize}

\begin{tcolorbox}[breakable,colback=teal!4,colframe=teal!35!black,title={Before you move on}]
What does reclamare mean? (\textquotedblleft{}To complain, to protest; to claim\textquotedblright{}.) Say it once more.
\end{tcolorbox}
\section[la lamentela]{la lamentela — the grumble, the complaint}

\label{lesson:IT-C438-lamentela}



\begin{quote}
Before the new one: say the Italian for the complaint, then the Italian for to complain.
\end{quote}

\subsection*{You'll want to know: la lamentela}
\textbf{la lamentela} — \textquotedblleft{}the grumble, the complaint\textquotedblright{}.

\begin{quote}
Ricevono molte lamentele per la colazione.
\end{quote}

\begin{quote}
\emph{They get a lot of complaints about breakfast.}
\end{quote}

Less formal than \emph{il reclamo}; \emph{lamentarsi} — to complain.

\begin{cousinweb}
\emph{lamentela} is built on \emph{lamentare} 'to lament', which is inherited from Latin \emph{lāmentārī} (EtymDB), from \emph{lāmentum} 'wailing' (Lewis \& Short).
\end{cousinweb}

\begin{grammarlens}[title={What you've built}]
One more word for telling, explaining and describing at B1.
\end{grammarlens}

\subsection*{Your turn}
\begin{itemize}
\raggedright
  \item \emph{Say it:} \emph{la lamentela}
  \item \emph{Say it:} \emph{la lamentela}, once more
  \item \emph{Say it:} a sentence of your own with \emph{la lamentela}
  \item \emph{Recall:} say the Italian for the complaint, then the Italian for to complain, then say \emph{la lamentela} again
\end{itemize}

\begin{tcolorbox}[breakable,colback=teal!4,colframe=teal!35!black,title={Before you move on}]
What does la lamentela mean? (\textquotedblleft{}The grumble, the complaint\textquotedblright{}.) Say it once more.
\end{tcolorbox}
\section[rimediare]{rimediare — to put right, to make up for}

\label{lesson:IT-C438-rimediare}



\begin{quote}
Before the new one: say the Italian for to complain, then the Italian for the grumble.
\end{quote}

\subsection*{You'll want to know: rimediare}
\textbf{rimediare} — \textquotedblleft{}to put right, to make up for\textquotedblright{}.

\begin{quote}
L'albergo ha rimediato con una camera più grande.
\end{quote}

\begin{quote}
\emph{The hotel made up for it with a bigger room.}
\end{quote}

\emph{Il rimedio} — the remedy.

\begin{cousinweb}
\emph{rimediare} is built on \emph{rimedio} 'remedy', a learned borrowing from Latin \emph{remedium} 'cure', from \emph{medērī} 'to heal' (Lewis \& Short).
\end{cousinweb}

\begin{grammarlens}[title={What you've built}]
One more word for telling, explaining and describing at B1.
\end{grammarlens}

\subsection*{Your turn}
\begin{itemize}
\raggedright
  \item \emph{Say it:} \emph{rimediare}
  \item \emph{Say it:} \emph{rimediare}, once more
  \item \emph{Say it:} a sentence of your own with \emph{rimediare}
  \item \emph{Recall:} say the Italian for to complain, then the Italian for the grumble, then say \emph{rimediare} again
\end{itemize}

\begin{tcolorbox}[breakable,colback=teal!4,colframe=teal!35!black,title={Before you move on}]
What does rimediare mean? (\textquotedblleft{}To put right, to make up for\textquotedblright{}.) Say it once more.
\end{tcolorbox}
\section[risarcire]{risarcire — to compensate, to pay damages}

\label{lesson:IT-C438-risarcire}



\begin{quote}
Before the new one: say the Italian for the grumble, then the Italian for to put right.
\end{quote}

\subsection*{You'll want to know: risarcire}
\textbf{risarcire} — \textquotedblleft{}to compensate, to pay damages\textquotedblright{}.

\begin{quote}
La compagnia deve risarcire i passeggeri.
\end{quote}

\begin{quote}
\emph{The company must compensate the passengers.}
\end{quote}

\emph{Il risarcimento} — compensation. Conjugated like \emph{finire}.

\begin{cousinweb}
\emph{risarcire} is inherited from Latin \emph{resarcīre} 'to patch or mend again' (EtymDB), from \emph{sarcīre} 'to mend'.
\end{cousinweb}

\begin{grammarlens}[title={What you've built}]
One more word for telling, explaining and describing at B1.
\end{grammarlens}

\subsection*{Your turn}
\begin{itemize}
\raggedright
  \item \emph{Say it:} \emph{risarcire}
  \item \emph{Say it:} \emph{risarcire}, once more
  \item \emph{Say it:} a sentence of your own with \emph{risarcire}
  \item \emph{Recall:} say the Italian for the grumble, then the Italian for to put right, then say \emph{risarcire} again
\end{itemize}

\begin{tcolorbox}[breakable,colback=teal!4,colframe=teal!35!black,title={Before you move on}]
What does risarcire mean? (\textquotedblleft{}To compensate, to pay damages\textquotedblright{}.) Say it once more.
\end{tcolorbox}
